% Part: lambda-calculus
% Chapter: introduction
% Section: syntax

\documentclass[../../../include/open-logic-section]{subfiles}

\begin{document}

\olfileid{lam}{int}{syn}
\olsection{Sintaksis Kalkulus Lambda}

Kita miwiti nganggo rerangken variabel $x$, $y$,
$z$,~\dots lan sawetara simbol konstanta $a$, $b$,
$c$,~\dots. Himpunan suku ditetepake kanthi induktif
kaya mangkene:
\begin{enumerate}
\item Saben variabel iku suku.
\item Saben konstanta iku suku.
\item Yen $M$ lan $N$ iku suku, mula $(MN)$ uga suku.
\item Yen $M$ iku suku lan $x$ iku variabel, mula
  $(\lambd[x][M])$ uga suku.
\end{enumerate}
Suku kang wujude $(MN)$ diarani \emph{aplikasi},
dene suku kang wujude $(\lambd[x][M])$
diarani \emph{abstraksi}.

Sistem kang ora ngemot konstanta babar pisan diarani
kalkulus lambda \emph{murni}. Kita mligine bakal
nggarap kalkulus-$\lambd$ murni, mula kabeh aksara
cilik makili variabel. Aksara gedhe ($M$, $N$,
lan sateruse) digunakake kanggo makili suku
kalkulus-$\lambd$.

Kita bakal netepi sawetara paugeran notasi:

\begin{conv}
\begin{enumerate}
\item Yen tandha kurung ora ditulis, aplikasi
  ditindakake saka kiwa menyang tengen.
  Umpamane, yen $M$, $N$, $P$, lan $Q$ iku
  suku, mula $MNPQ$ iku singkatan saka
  $(((MN)P)Q)$.
\item Maneh, yen tandha kurung ora ditulis,
  abstraksi lambda diwenehi cakupan paling
  amba. Umpamane, $\lambd[x][MNP]$
  diwaca $(\lambd[x][((MN)P)])$.
\item Lambda bisa digunakake kanggo
  ngabstraksi sawetara variabel bebarengan.
  Umpamane, $\lambd[xyz][M]$ iku
  singkatan saka
  $\lambd[x][\lambd[y][\lambd[z][M]]]$.
\end{enumerate}
\end{conv}

Umpamane,
\[
\lambd[xy][xxyx \lambd[z][xz]]
\]
iku singkatan saka
\[
\lambd[x][\lambd[y][((((xx)y)x)(\lambd[z][(xz)]))]].
\]
Paugeran iki kudu dielingi. Wiwitane bisa nggawe
mumet, nanging mengko bakal kulina,
lan suwene-suwene paugeran iki ora semumet
tinimbang kudu ngurusi kurung kang ruwet
ora karuan.

Rong suku kang mung beda jeneng variabel kaiket
diarani ekuivalen-$\alpha$; umpamane,
$\lambd[x][x]$ lan $\lambd[y][y]$.
Luwih trep yen loro iku dianggep minangka suku
kang ``padha''; kanthi tembung liya,
nalika kita nyebut $M$ lan $N$ padha, maksude
uga ``nganti owah-owahan jeneng variabel
kaiket''. Muncule variabel kang diiket
dening sawijining $\lambd$ ing cakupane diarani
``kaiket'', dene kemunculan liyane
diarani ``bebas''.
Ora ana variabel bebas ing tuladha sadurunge;
nanging ing
\[
(\lambd[z][yz])x
\]
$y$ lan $x$ iku bebas, dene $z$ kaiket.

\end{document}
